\chapter{Mission 2 and 3: Move It, Catch It}
\label{ch:kids2}

Your square sits still. A game where nothing moves is not much of a game. In
this chapter you make the square obey the buttons, and then you build a whole
game called \textbf{Lemon Catcher}.

\section{The game pad}

Under the game screen there is a game pad. It has a cross with four arrows, two
round buttons called \kbtn{A} and \kbtn{B}, and two small buttons called
\kbtn{SELECT} and \kbtn{START}. You can click them with the mouse or touch them
with a finger. The arrow keys on the keyboard work too, with \kbtn{Z} for A and
\kbtn{X} for B.

\section{If}

The red block in the \textbf{Input \& Logic} drawer asks a question:

\kinput{if button LEFT pressed}

The blocks you put \emph{inside} it run only when the answer is yes. If the
button is not pressed, the computer skips them.

\begin{prgconcept}
\textbf{If} lets a program choose. \emph{If} the button is pressed, do this.
Otherwise, do nothing and carry on.
\end{prgconcept}

\section{Change}

The purple block \textbf{change} adds a number to a variable.

\kstate{change player\_x by 3}

If player\_x was 152, it becomes 155. Next frame, 158. Then 161. The square
slides to the right.

To go left, x must get \emph{smaller}. So we add a number below zero:

\kstate{change player\_x by -3}

\begin{prgmission}[Mission 2, step 1: move left and right]
Open your First Square project. Drag \textbf{every frame update} into the
workspace and build this:

\kgame{every frame update}
\kinput[6mm]{if button LEFT pressed}
\kstate[12mm]{change player\_x by -3}
\kinput[6mm]{if button RIGHT pressed}
\kstate[12mm]{change player\_x by 3}

Press \kbtn{Save} and \kbtn{Play}. Use the arrows!
\end{prgmission}

\begin{prgtry}
\begin{enumerate}
  \item Change 3 to 1. Then to 10. What does this number do?
  \item Add two more \textbf{if} blocks for UP and DOWN. They must change
        \textbf{player\_y}. Which one needs the minus sign?
\end{enumerate}
\end{prgtry}

\section{Do not fall off the world}

Hold the right arrow for a long time. The square walks off the screen and never
comes back! The variable keeps growing: 320, 330, 340\dots

The purple block \textbf{keep} is a fence. It stops a number from going too low
or too high.

\kstate{keep player\_x between 0 and 304}

Why 304 and not 320? The square is 16 wide. Its left side is at x. If x is 304,
its right side is at $304 + 16 = 320$, exactly the edge.

\begin{prgmission}[Mission 2, step 2: the fence]
Add two \textbf{keep} blocks at the bottom of \emph{every frame update}:

\kstate[6mm]{keep player\_x between 0 and 304}
\kstate[6mm]{keep player\_y between 0 and 184}

Play again. Can the square escape now?
\end{prgmission}

\section{Watching the numbers}

Under the game pad there is a grey line. It shows your variables and their
numbers while the game runs.

Press \kbtn{Play} once more to pause. Now press \kbtn{Step}. The game moves
forward by \emph{one} frame only. Hold the right arrow on the game pad and press
\kbtn{Step} again and again. Watch player\_x: 152, 155, 158\dots

\begin{prgconcept}
\textbf{Step} runs one frame and stops. It lets you slow the computer down
until you can see exactly what your blocks do.
\end{prgconcept}

\section{Mission 3: Lemon Catcher}

Time for a real game. Lemons fall from the sky. You move a basket to catch
them. A green bar grows when you catch one. A red bar grows when you miss.

Start a \textbf{new project} called \textbf{Lemon Catcher}.

\subsection*{The boxes we need}

\begin{center}
\begin{tabular}{@{}ll@{}}
\toprule
Variable & What it remembers \\
\midrule
basket\_x & where the basket is, left to right \\
lemon\_x & where the lemon is, left to right \\
lemon\_y & how far the lemon has fallen \\
score & how many lemons you caught \\
misses & how many lemons you missed \\
\bottomrule
\end{tabular}
\end{center}

\begin{prgmission}[Mission 3, step 1: get ready]
\kgame{when game starts}
\kstate[6mm]{set basket\_x to 144}
\kstate[6mm]{set lemon\_x to 150}
\kstate[6mm]{set lemon\_y to 0}
\kstate[6mm]{set score to 0}
\kstate[6mm]{set misses to 0}
\end{prgmission}

\begin{prgmission}[Mission 3, step 2: draw everything]
\kgame{every frame draw}
\kdraw[6mm]{clear screen BLACK}
\kdraw[6mm]{draw rectangle x lemon\_x y lemon\_y w 10 h 10 color YELLOW}
\kdraw[6mm]{draw rectangle x basket\_x y 180 w 32 h 8 color ORANGE}
\kdraw[6mm]{draw rectangle x 0 y 196 w 320 h 4 color BLUE}

Press \kbtn{Play}. You see a lemon at the top, a basket, and a blue floor.
Nothing moves yet. That is fine. We add one thing at a time.
\end{prgmission}

\begin{prgmission}[Mission 3, step 3: move the basket, drop the lemon]
\kgame{every frame update}
\kinput[6mm]{if button LEFT pressed}
\kstate[12mm]{change basket\_x by -4}
\kinput[6mm]{if button RIGHT pressed}
\kstate[12mm]{change basket\_x by 4}
\kstate[6mm]{keep basket\_x between 0 and 288}
\kstate[6mm]{change lemon\_y by 2}

Play. The basket moves. The lemon falls\dots and falls through the floor.
\end{prgmission}

\section{Touching}

How does the game know the lemon is in the basket? With the other red block:

\kinput{if rect \dots{} touches rect \dots}

``Rect'' is short for \emph{rectangle}. The block has eight numbers: x, y, w
and h for the first rectangle, then x, y, w and h for the second. If the two
rectangles overlap, even a little, the blocks inside run.

\begin{center}
\begin{tikzpicture}[font=\sffamily\small]
  \draw[fill=prgamber!70,draw=prgamber!60!black] (0,0.4) rectangle (0.8,1.2);
  \draw[fill=prgred!50,draw=prgred] (1.6,0) rectangle (4.0,0.5);
  \node at (2,-0.5){not touching};
  \draw[fill=prgred!50,draw=prgred] (6.6,0) rectangle (9.0,0.5);
  \draw[fill=prgamber!70,draw=prgamber!60!black] (7.2,0.3) rectangle (8.0,1.1);
  \node at (7.8,-0.5){touching!};
\end{tikzpicture}
\captionof{figure}{Two rectangles touch when they overlap.}
\end{center}

When you catch a lemon, three things should happen: you get a point, and a new
lemon appears at the top, somewhere new. For ``somewhere new'' there is a
purple block that rolls a dice:

\kstate{set lemon\_x to random from 0 to 310}

\begin{prgmission}[Mission 3, step 4: catch!]
Add this at the bottom of \emph{every frame update}. Type the eight values
carefully: the lemon first, then the basket.

\kinput[6mm]{if rect lemon\_x lemon\_y 10 10 touches rect basket\_x 180 32 8}
\kstate[12mm]{change score by 1}
\kstate[12mm]{set lemon\_y to 0}
\kstate[12mm]{set lemon\_x to random from 0 to 310}
\end{prgmission}

\begin{prgmission}[Mission 3, step 5: miss!]
The blue floor is a rectangle too. If the lemon touches the floor, you missed.

\kinput[6mm]{if rect lemon\_x lemon\_y 10 10 touches rect 0 196 320 4}
\kstate[12mm]{change misses by 1}
\kstate[12mm]{set lemon\_y to 0}
\kstate[12mm]{set lemon\_x to random from 0 to 310}
\end{prgmission}

\section{Showing the score}

The player cannot see inside your variables. Let us draw the score as a bar
that grows. The trick: a width does not have to be a plain number. It can be a
little sum.

\begin{prgmission}[Mission 3, step 6: score bars]
Add these to \emph{every frame draw}, after \textbf{clear screen}:

\kdraw[6mm]{draw text LEMONS at x 8 y 4 fg GREEN bg BLACK}
\kdraw[6mm]{draw rectangle x 64 y 4 w score * 3 h 8 color GREEN}
\kdraw[6mm]{draw text MISSED at x 8 y 16 fg RED bg BLACK}
\kdraw[6mm]{draw rectangle x 64 y 16 w misses * 3 h 8 color RED}

The star \textbf{*} means ``times''. With 5 lemons the bar is $5 \times 3 = 15$
pixels wide. Play, and watch the bars race!
\end{prgmission}

\begin{prgtry}
Make the game your own.
\begin{enumerate}
  \item Make the lemon fall faster. Which number is it?
  \item Make the basket smaller. You must change it in \emph{two} places. Why?
  \item Add a second fruit with its own variables: orange\_x and orange\_y.
  \item Make a ``bad'' fruit that takes away a point when you catch it.
\end{enumerate}
\end{prgtry}

\begin{prgwarn}
If your basket catches nothing, check the \textbf{touches} block. The numbers
there must be the same as the numbers in your \textbf{draw rectangle} blocks.
The computer does not look at the picture. It only looks at the numbers.
\end{prgwarn}

\begin{prgcheck}
You finished Missions~2 and~3 if you can:
\begin{itemize}
  \item explain what \textbf{if} does;
  \item make something move with a button;
  \item say why the fence stops at 304 and not at 320;
  \item use \kbtn{Step} to watch a variable change;
  \item explain how the game knows a lemon was caught.
\end{itemize}
\end{prgcheck}

\begin{prggrownup}
Reference projects: \fname{moving\_square.blocks.json} and
\fname{lemon\_catcher.blocks.json} in \fname{examples/blocks\_first\_steps}.
The Blocks language deliberately has no general comparison block; the ``floor
is a rectangle'' idea in step~5 is how a threshold test is expressed with the
collision block. A complete, playable C rewrite with a numeric score and a
game-over state is in Appendix~\ref{app:worked}, for older siblings and for the
C track.
\end{prggrownup}
